\chapter{Mathematical and Geometric Algorithms}
\label{ch:math}

Everything the camera sees has now become 21 points with names (Chapter~\ref{ch:engine}). This
chapter is the part of the project that is \emph{ours}: turning those points into answers. It is the
only chapter where the reasoning is original rather than borrowed from a library, and it is where
every interesting bug lives.

The chapter moves from the simplest question to the most subtle:

\begin{enumerate}
  \item Which way is ``up'' on a screen, and why does that invert the finger test?
        (\S\ref{sec:screencs})
  \item Is the thumb out? --- the rule that survives rotation, and the physical pose that still
        defeats it. (\S\ref{sec:thumb})
  \item Which side of the hand is facing the camera? (\S\ref{sec:palm})
  \item How do we stop the numbers shaking? (\S\ref{sec:ema})
  \item What gesture is this, and how do we stop the label flickering? (\S\ref{sec:gesture})
\end{enumerate}

\section{The screen coordinate system, and the inverted y-axis}
\label{sec:screencs}

\begin{mathdeepdive}{Two coordinate systems, one picture}
In the mathematics you learned at school, the \textbf{Cartesian plane} has its origin in the middle
or at the bottom-left, the $x$-axis grows to the right, and \textbf{$y$ grows upward}. A point
``higher up'' has a \emph{larger} $y$.

A computer screen inherited a different convention from the cathode-ray tubes of the 1960s: the
electron beam scanned from the top-left corner, so that became the origin, and the beam's downward
sweep became the positive $y$ direction. **On a screen, $y$ grows downward.** A point higher up has a
\emph{smaller} $y$.

The two systems are related by a vertical flip:

$$y_{\text{cartesian}} = H - y_{\text{screen}}$$

where $H$ is the image height in pixels. Nothing else changes; $x$ grows to the right in both.
\end{mathdeepdive}

\begin{figure}[h]
\centering
\begin{tikzpicture}[every node/.style={font=\small\sffamily}]
  % ---- Cartesian (left) ----
  \begin{scope}[shift={(0,0)}]
    \draw[slateline, line width=0.8pt] (-0.4,-0.4) rectangle (3.6,3.0);
    \draw[-{Stealth[length=2mm]}, navy, line width=1pt] (0,0) -- (3.3,0)
          node[right, text=navy] {$x$};
    \draw[-{Stealth[length=2mm]}, navy, line width=1pt] (0,0) -- (0,2.7)
          node[above, text=navy] {$y$};
    \fill[redx] (2.0,1.8) circle (2.2pt);
    \node[right=2pt, text=redx] at (2.0,1.8) {$P$};
    \draw[cyanx, dashed] (0,1.8) -- (2.0,1.8);
    \draw[cyanx, dashed] (2.0,0) -- (2.0,1.8);
    \node[below, text=cyanx] at (1.0,0.05) {$x$};
    \node[left, text=cyanx] at (0.05,0.9) {$y$};
    \node[below=0.35cm, text=navy, font=\small\sffamily\bfseries] at (1.6,-0.5)
         {Cartesian: $y$ grows \emph{up}};
  \end{scope}

  % ---- Screen (right) ----
  \begin{scope}[shift={(5.4,0)}]
    \draw[slateline, line width=0.8pt, fill=slatebg] (-0.4,-0.4) rectangle (3.6,3.0);
    \draw[-{Stealth[length=2mm]}, navy, line width=1pt] (-0.4,-0.4) -- (3.3,-0.4)
          node[right, text=navy] {$x$};
    \draw[-{Stealth[length=2mm]}, navy, line width=1pt] (-0.4,-0.4) -- (-0.4,2.7);
    \node[above, text=navy] at (-0.4,2.7) {$y$};
    \node[below left, text=navy, font=\scriptsize] at (-0.4,-0.4) {$(0,0)$};
    \fill[redx] (2.0,1.3) circle (2.2pt);
    \node[right=2pt, text=redx] at (2.0,1.3) {$P$};
    \draw[cyanx, dashed] (-0.4,1.3) -- (2.0,1.3);
    \draw[cyanx, dashed] (2.0,-0.4) -- (2.0,1.3);
    \node[left, text=cyanx] at (0.75,0.85) {$y$};
    \node[below, text=cyanx] at (2.0,-0.15) {$\dots$};
    \node[below=0.35cm, text=navy, font=\small\sffamily\bfseries] at (1.6,-0.5)
         {Screen: $y$ grows \emph{down}};
    \node[right=0.1cm, text=slatetext, font=\scriptsize, align=left] at (3.4,1.3)
         {origin\\top-left};
  \end{scope}

  % relationship
  \draw[-{Stealth[length=2mm]}, tealx, line width=1pt] (4.0,1.3) -- (4.9,1.3);
  \node[text=tealx, font=\small\sffamily, above] at (4.45,1.35) {$H-y$};
\end{tikzpicture}
\caption{The same point $P$ in the two conventions. In the Cartesian frame the vertical arrow
points up and $P$ has a positive $y$. On a screen the vertical arrow points down and the identical
point now has $y$ measured from the top-left corner. Every ``my marker is mirrored vertically'' bug
is this figure being forgotten.}
\label{fig:coords}
\end{figure}

\subsection{Why a fingertip being ``above'' its knuckle means a smaller $y$}

This single fact is the whole basis of the four-finger test.

\begin{mathdeepdive}{Deriving the finger test}
Hold your hand up in front of the camera with the index finger extended. In the image, the
fingertip is physically higher than the knuckle.

Because $y$ grows \emph{downward}, ``physically higher'' means ``numerically smaller'':

$$y_{\text{tip}} < y_{\text{knuckle}} \quad \Longleftrightarrow \quad \text{the fingertip is higher on screen}$$

Now bend the finger so the tip curls down toward the palm. The tip drops below the knuckle in the
image, so its $y$ becomes \emph{larger}:

$$y_{\text{tip}} > y_{\text{knuckle}} \quad \Longleftrightarrow \quad \text{the finger is folded}$$

The test is therefore:

$$\text{finger extended} \iff y_{\text{tip}} < y_{\text{pip}}$$

and the comparison is \emph{strict} --- an exactly equal $y$ counts as folded, which costs nothing
and keeps the rule total (it always returns a definite answer).

\textbf{Worked example.} Say the index fingertip is at pixel $(346, 235)$ and the index PIP joint at
$(350, 280)$. Then $y_{\text{tip}} = 235 < y_{\text{pip}} = 280$, so the finger is extended --- and
indeed $45$ pixels of difference is a clearly raised finger. If instead the tip were at
$(352, 305)$, then $305 > 280$ and the finger is folded.
\end{mathdeepdive}

The four long fingers are tested with four (tip, knuckle) index pairs from \code{hand_logic.py}:

\begin{codecard}
\begin{lstlisting}
FINGER_PAIRS = [(8, 6), (12, 10), (16, 14), (20, 18)]
\end{lstlisting}
\end{codecard}

\begin{center}
\rowcolors{2}{slatebg}{white}
\begin{tabular}{@{}l l l l@{}}
\toprule
\rowcolor{white}
\textbf{Finger} & \textbf{Tip index} & \textbf{Knuckle (PIP) index} & \textbf{Test} \\
\midrule
Index  & 8  & 6  & $y_8 < y_6$ \\
Middle & 12 & 10 & $y_{12} < y_{10}$ \\
Ring   & 16 & 14 & $y_{16} < y_{14}$ \\
Pinky  & 20 & 18 & $y_{20} < y_{18}$ \\
\bottomrule
\end{tabular}
\end{center}

\begin{edgetrap}{The y-rule's hidden assumption: fingers must point roughly up}
The test compares $y$ only. That is a comparison of \emph{vertical} position, and it is only
meaningful while the fingers point roughly toward the top of the image.

Point your index finger horizontally to the right (thumb up, like a finger-gun) and the tip and
knuckle have almost the same $y$. The tip may even be slightly lower. The rule will report the finger
as folded while it is plainly extended.

This is a genuine limitation, not a bug, and it is why the tracker is documented as a
\emph{upright-hand} tool. Making it fully rotation-invariant would require comparing the fingertip
against the palm axis --- computing the finger's direction relative to the wrist-to-middle-knuckle
vector and asking whether the tip is farther along that axis than the knuckle. That is a real
technique, and the thumb rule below uses exactly that family of idea (a distance from a fixed palm
reference), but applying it to all four fingers is more machinery than this project needs.

\textbf{Practical consequence for you as a user:} keep your hand upright. \textbf{Practical
consequence as an engineer:} if you extend this project, this is the first rule to replace.
\end{edgetrap}

\section{The thumb, and the search for invariance}
\label{sec:thumb}

The thumb is the hard one. It moves on a different axis from the other four fingers, it is
foreshortened by perspective when it points at the camera, and its apparent left-right position
depends on which hand it is \emph{and} which way the palm is facing.

\subsection{The obvious rule, and exactly how it breaks}

The original specification proposed a horizontal comparison. Because a webcam feed is mirrored, and
because MediaPipe's handedness label is predicted under a mirrored-feed assumption, the rule was
written as:

$$\text{Right} \Rightarrow \text{thumb extended if } x_{\text{tip}} < x_{\text{ip}}$$
$$\text{Left} \Rightarrow \text{thumb extended if } x_{\text{tip}} > x_{\text{ip}}$$

\begin{edgetrap}{Why the horizontal rule fails: palm inversion}
The rule silently encodes a physical assumption --- \emph{the palm faces the camera}. Turn your hand
around so the back faces the camera, and everything the rule knows becomes false.

Concretely: hold up your right hand, palm toward the camera, thumb out. The thumb points to one side
of the image. Now rotate your wrist 180\textdegree\ so the back of the hand faces the camera, keeping
the thumb extended. The thumb is now on the \emph{other} side of the image, while the hand is still
your right hand and the label still reads ``Right''. The comparison $x_{\text{tip}} < x_{\text{ip}}$
now returns the opposite answer.

The result: the finger count flickers as you rotate your hand, and the gesture classification
collapses with it. Rotation of the wrist alone --- no change in pose at all --- flips the answer.

There is a second, subtler failure: the comparison is purely horizontal, so a hand rotated in the
image plane (say, 45\textdegree) has a thumb whose $x_{\text{tip}} - x_{\text{ip}}$ difference is
scaled by $\cos 45^\circ \approx 0.71$, shrinking the margin and making the test fragile near the
decision boundary.
\end{edgetrap}

The rule is kept in the code because the specification asked for it, reachable as
\code{--thumb-mode x}, but it is not the default.

\subsection{The distance rule}

\begin{keyidea}
\textbf{The idea in one sentence.}
However the hand is rotated, an \textbf{extended thumb reaches away from the palm}, while a folded
thumb is tucked \emph{across} the palm --- so measure how far the thumb tip and the thumb joint sit
from a fixed palm reference, and see which is farther.
\end{keyidea}

The fixed reference is landmark 17, the pinky MCP --- the far corner of the palm, the point the thumb
moves \emph{toward} when it folds and \emph{away} from when it extends.

\begin{mathdeepdive}{The rule, derived and justified}
$$\text{thumb extended} \iff d\big(P_4,\ P_{17}\big) > d\big(P_3,\ P_{17}\big)$$

where $P_4$ is the thumb tip, $P_3$ is the thumb IP joint, $P_{17}$ is the pinky MCP, and $d$ is the
ordinary Euclidean distance

$$d(A,B) = \sqrt{(A_x - B_x)^2 + (A_y - B_y)^2}.$$

\textbf{Why a distance and not a coordinate.} A coordinate is tied to the frame of reference: $x$
means ``how far right on the screen'', which changes when the hand rotates. A distance is an
\emph{intrinsic} quantity --- it depends only on the two points themselves.

\textbf{The invariance argument, stated precisely.} Suppose the hand is rotated in the image plane by
an angle $\theta$ about any centre $C$. The rotation matrix is

$$R_\theta = \begin{pmatrix} \cos\theta & -\sin\theta \\ \sin\theta & \cos\theta \end{pmatrix},
\qquad R_\theta^{\mathsf T} R_\theta = I.$$

Every landmark is transformed as $P' = C + R_\theta (P - C)$. Therefore the difference
$P_4' - P_{17}' = R_\theta (P_4 - P_{17})$, and

\begin{align*}
d(P_4', P_{17}') &= \big\| R_\theta (P_4 - P_{17}) \big\| \\
                 &= \sqrt{(P_4-P_{17})^{\mathsf T} R_\theta^{\mathsf T} R_\theta (P_4-P_{17})} \\
                 &= \sqrt{(P_4-P_{17})^{\mathsf T} (P_4-P_{17})} = d(P_4, P_{17}).
\end{align*}

The distance is \emph{exactly} unchanged by any rotation, and the same holds for $d(P_3, P_{17})$.
Since both sides of the comparison are preserved, the comparison's result cannot change. The edge
case is equality --- the two distances are equal exactly when the thumb tip and IP joint lie on a
circle centred on the pinky knuckle, which happens only in a degenerate straight-thumb pose; the
strict \code{>} then reports ``folded'', which is the safe answer for a thumb that is not clearly
away from the palm.

\textbf{Handedness-independence.} Nothing in the expression refers to ``left'' or ``right''. It is a
statement about three points, so it cannot be affected by a label that may itself be wrong.
\end{mathdeepdive}

\begin{mathdeepdive}{Worked example, with numbers}
This is the actual synthetic hand from the project's test suite (a right hand, palm to camera,
fingers up, mirrored view).

\textbf{Thumb extended}
\begin{align*}
P_4 &= (430, 320) \quad \text{(thumb tip)} \\
P_3 &= (415, 345) \quad \text{(thumb IP)} \\
P_{17} &= (262, 335) \quad \text{(pinky MCP)}
\end{align*}

\begin{align*}
d(P_4, P_{17}) &= \sqrt{(430-262)^2 + (320-335)^2} \\
               &= \sqrt{168^2 + (-15)^2} = \sqrt{28224 + 225} = \sqrt{28449} \approx 168.7
\end{align*}

\begin{align*}
d(P_3, P_{17}) &= \sqrt{(415-262)^2 + (345-335)^2} \\
               &= \sqrt{153^2 + 10^2} = \sqrt{23409 + 100} = \sqrt{23509} \approx 153.3
\end{align*}

Since $168.7 > 153.3$, the thumb is reported \textbf{extended}. \checkmark

\textbf{Thumb folded.} Now the tip folds across the palm to $P_4 = (375, 350)$, with the IP joint at
$P_3 = (395, 355)$:

\begin{align*}
d(P_4, P_{17}) &= \sqrt{113^2 + 15^2} = \sqrt{12769+225} = \sqrt{12994} \approx 114.0 \\
d(P_3, P_{17}) &= \sqrt{133^2 + 20^2} = \sqrt{17689+400} = \sqrt{18089} \approx 134.5
\end{align*}

Now $114.0 < 134.5$, so the thumb is reported \textbf{folded}. \checkmark

Notice the sizes involved: the margins are roughly $15$ pixels and $20$ pixels respectively. That is
a comfortable margin at $640 \times 480$ --- large enough that landmark jitter of one or two pixels
cannot flip the answer.
\end{mathdeepdive}

\subsection{The knife hand: when the distance rule lies}

\begin{edgetrap}{The adducted-thumb false positive}
Hold your hand flat, fingers straight, thumb pressed tightly against the side of the index finger ---
the ``karate chop'' or ``knife hand''.

The thumb tip is physically at the far end of the hand, near the index fingertip, while the thumb IP
joint is back near the palm. Distance to the pinky knuckle (landmark 17): the tip really \emph{is}
farther away than the IP joint. So the rule reports the thumb as \textbf{extended}, and the finger
count is one too high.

This is not a coding error; it is the rule doing exactly what it says. The failure is that the rule
only ever asks ``is the thumb away from the pinky side of the palm?'', and in a knife hand the answer
is genuinely yes.

\textbf{The candidate fix} --- a secondary clearance gate against the index MCP (landmark 5), which
the thumb must also be distant from:

$$\text{dist}(P_4, P_5) > 1.2 \cdot \text{dist}(P_3, P_5)$$

The factor $1.2$ encodes ``the tip must be \emph{clearly} farther from the index knuckle than the IP
joint, not merely farther''. In a knife hand the tip sits close to the index finger, so this second
test fails and the false positive is suppressed.

\textbf{Why this project does not enable it by default.} The gate trades one error for another. On a
thumb pointing sideways across the image --- a genuine ``thumbs up'' rotated 90\textdegree\ by wrist
tilt --- the tip can pass close to the index knuckle and the gate would suppress a \emph{correct}
extended result. The factor $1.2$ is a tuned constant with no derivation, and it varies with hand
size and camera distance because it compares raw pixel distances.

\textbf{The engineering position taken here:} ship the rule that is rotation-invariant and
documented, document the pose that defeats it, and leave the tuning to whoever has the hardware in
front of them. The limitation is stated plainly in the project's README rather than hidden.
\end{edgetrap}

\subsection{Keeping the old rule reachable, and what happens with no label}

The code implements both rules, with the distance rule as the default and a defined behaviour when
the handedness label is unavailable:

\begin{codecard}
\begin{lstlisting}
if thumb_mode == THUMB_X and hand_label in ("Left", "Right"):
    tip, ip = pixel_landmarks[THUMB_TIP], pixel_landmarks[THUMB_IP]
    thumb = tip[0] < ip[0] if hand_label == "Right" else tip[0] > ip[0]
else:
    thumb = _distance(pixel_landmarks[THUMB_TIP], pixel_landmarks[PINKY_MCP]) > _distance(
        pixel_landmarks[THUMB_IP], pixel_landmarks[PINKY_MCP]
    )
\end{lstlisting}
\end{codecard}

\begin{itemize}
  \item The condition is a conjunction: the naive rule needs \emph{both} \code{thumb\_mode == "x"}
        \emph{and} a usable label.
  \item If the label is \code{"Unknown"} --- which happens when MediaPipe's handedness list is
        missing or short --- the \code{else} branch applies \textbf{the distance rule}, not
        \code{False}. A missing label degrades to a working rule rather than silently losing the
        thumb.
  \item \code{tip[0]} is the $x$ component of the tip tuple; \code{tip[1]} would be $y$. Keeping
        coordinate access as indexed tuples rather than named objects is what makes this expression
        short.
\end{itemize}

\section{The palm normal: which side of the hand am I seeing?}
\label{sec:palm}

Finger counting does not need this. It is included because it is exactly the kind of diagnostic
telemetry the specification asks for, and because it makes the thumb rule's invariance visible while
the program runs.

\begin{mathdeepdive}{Two vectors and a cross product}
Take two vectors from the wrist $P_0$:

\begin{align*}
\vec{v}_1 &= P_5 - P_0 \quad \text{(wrist} \to \text{index knuckle)} \\
\vec{v}_2 &= P_{17} - P_0 \quad \text{(wrist} \to \text{pinky knuckle)}
\end{align*}

In 3D, the palm's normal vector is their cross product $\vec{v}_1 \times \vec{v}_2$. Writing
$\vec{v}_1 = (a_x, a_y, a_z)$ and $\vec{v}_2 = (b_x, b_y, b_z)$:

$$\vec{v}_1 \times \vec{v}_2 =
\begin{pmatrix} a_y b_z - a_z b_y \\ a_z b_x - a_x b_z \\ a_x b_y - a_y b_x \end{pmatrix}$$

The project uses \textbf{only the $z$ component}, because the landmarks' $x$ and $y$ are pixel
coordinates while $z$ is relative depth --- mixing them would be meaningless. With
$a = (x_5 - x_0,\ y_5 - y_0)$ and $b = (x_{17} - x_0,\ y_{17} - y_0)$:

\begin{align*}
Z_{\text{normal}} &= a_x b_y - a_y b_x \\
 &= (x_5 - x_0)(y_{17} - y_0) - (y_5 - y_0)(x_{17} - x_0)
\end{align*}

This is the signed area (times two) of the triangle wrist--index-knuckle--pinky-knuckle. Its
\textbf{sign} records the rotational order of the two vectors --- which way you turn going from
$\vec{v}_1$ to $\vec{v}_2$ --- and the two sides of a hand produce opposite orders.

\textbf{Worked example.} Using the test suite's synthetic palm-facing hand
$P_0 = (320,420)$, $P_5 = (355,320)$, $P_{17} = (262,335)$:

\begin{align*}
Z &= (355-320)(335-420) - (320-420)(262-320) \\
  &= (35)(-85) - (-100)(-58) \\
  &= -2975 - 5800 = -8775
\end{align*}

A \emph{negative} $Z$ for a palm facing the camera. Mirror the same hand horizontally (which is
exactly what ``showing the back of your hand'' does to the knuckle ordering) and every $x$ difference
changes sign, so $Z$ becomes $+8775$ --- the sign flips, as it must.
\end{mathdeepdive}

\begin{figure}[h]
\centering
\begin{tikzpicture}[scale=1.25, every node/.style={font=\scriptsize\sffamily}]
  \handcoordinates
  \handbones

  % palm triangle
  \fill[tealx, opacity=0.12] (L0) -- (L5) -- (L17) -- cycle;
  \draw[tealx, dashed, line width=0.7pt] (L0) -- (L5) -- (L17) -- cycle;

  \foreach \i in {0,...,20} {\node[joint] at (L\i) {};}
  \foreach \i in {4,8,12,16,20} {\node[jointtip] at (L\i) {};}

  % vectors
  \draw[-{Stealth[length=2.6mm]}, cyanx, line width=1.6pt] (L0) -- (L5)
        node[midway, above left=-1pt, text=cyanx] {$\vec{v}_1$};
  \draw[-{Stealth[length=2.6mm]}, amberx, line width=1.6pt] (L0) -- (L17)
        node[midway, above right=-1pt, text=amberx] {$\vec{v}_2$};

  \node[right=0pt, text=navy, font=\scriptsize\sffamily\bfseries] at (L0) {$P_0$};
  \node[right=1pt, text=navy, font=\scriptsize\sffamily\bfseries] at (L5) {$P_5$};
  \node[left=1pt, text=navy, font=\scriptsize\sffamily\bfseries] at (L17) {$P_{17}$};

  % normal indicator: circle-with-dot is the convention for "out of the page"
  \draw[violetx, line width=1.1pt] (-2.10, 1.30) circle (0.20);
  \fill[violetx] (-2.10, 1.30) circle (0.06);
  \node[violetx, font=\scriptsize\sffamily, align=center] at (-2.10, 0.72)
       {$\vec{n}=\vec{v}_1\times\vec{v}_2$\\\emph{out of the screen}};
  \draw[violetx, line width=0.7pt, dashed] (-1.88, 1.30) -- (0.35, 1.30);

  % sign summary
  \node[draw=slateline, fill=slatebg, rounded corners=1mm, inner sep=3pt, align=left]
        at (1.3,-1.05)
        {$Z<0 \Rightarrow$ \textbf{Palm} \quad $Z>0 \Rightarrow$ \textbf{Back}};
\end{tikzpicture}
\caption{The palm normal. The wrist and the two outer knuckles define a plane; the sign of the
$z$-component of $\vec{v}_1 \times \vec{v}_2$ says whether that plane is presented palm-first or
back-first. Because the calculation uses only two differences, it is one subtraction cheaper than a
true 3D cross product and needs no depth data.}
\label{fig:palmnormal}
\end{figure}

\subsection{The sign convention, and how to remember it}

The sign-to-label mapping depends on the mirroring of the image, so the project states its convention
explicitly: \textbf{a mirrored (selfie) view, as produced by \code{cv2.flip(frame, 1)}}.

\begin{mathdeepdive}{Where the sign comes from}
Consider a right hand, palm to camera, fingers up, in a horizontally mirrored image.

\begin{itemize}
  \item In a mirrored view your right hand appears on the \emph{right} side of the picture, and its
        thumb --- pointing away from the palm --- appears on the \emph{far right} of the hand.
  \item Therefore the index knuckle $P_5$ lies to the \emph{right} of the wrist, and the pinky
        knuckle $P_{17}$ lies to the \emph{left}: $x_5 > x_0$ and $x_{17} < x_0$.
  \item Both knuckles are \emph{above} the wrist, and $y$ grows downward, so
        $y_5 < y_0$ and $y_{17} < y_0$.
\end{itemize}

Substituting these signs into the formula from the previous box,
$a_x b_y = (+)(-)$ and $a_y b_x = (-)(-)$, gives
$Z = (\text{negative}) - (\text{positive}) < 0$.

Now turn the hand over. Index and pinky knuckles swap sides, so $x_5 < x_0$ and $x_{17} > x_0$;
both $y$ relations are unchanged. Now $a_x b_y$ is negative and $a_y b_x$ is positive, giving
$Z > 0$.

Hence, for the mirrored convention: $Z<0 \Rightarrow$ \textbf{Palm facing}, $Z>0 \Rightarrow$
\textbf{Back facing}. The project encodes exactly this:

\begin{lstlisting}
def orientation_from_palm_z(palm_z):
    """Map the palm-normal sign to the HUD label ``Palm'' or ``Back''."""
    return "Palm" if palm_z < 0 else "Back"
\end{lstlisting}

The degenerate case $Z = 0$ (wrist, index knuckle and pinky knuckle collinear --- geometrically
impossible for a real hand, but reachable with synthetic or corrupted data) falls to \code{Back}
through the \code{else}. A test asserts this, so the behaviour is fixed rather than accidental.
\end{mathdeepdive}

\begin{edgetrap}{The convention is tied to the mirror, not to the hand}
Because the sign depends on the image being flipped, the label inverts in \code{--image} mode, where
a still photograph is deliberately \emph{not} mirrored. A photograph of a palm-first hand in
\code{--image} mode will report \code{Back}.

That is not a bug in the geometry; it is the geometry faithfully reporting an unmirrored image with a
mirrored-image convention. It is documented in the README, and it matters only for the diagnostic
label --- no finger decision depends on it.
\end{edgetrap}

\section{Smoothing: the exponential moving average}
\label{sec:ema}

\begin{mathdeepdive}{Sensor jitter, and why raw coordinates shake}
Even with a completely still hand, the reported landmark positions move by one to three pixels from
frame to frame. The causes are mundane and unavoidable:

\begin{itemize}
  \item camera sensor noise, especially in low light;
  \item micro-movements of your hand that you cannot consciously control;
  \item the model's own regression error, which varies slightly as the hand's appearance in the crop
        changes.
\end{itemize}

At 30 FPS, a dot drawn from raw coordinates therefore vibrates visibly --- it reads as ``cheap
program''. The fix is a filter: blend each new measurement with the previous smoothed value.

$$\mathbf{P}^{(t)}_{\text{smooth}} = \alpha\,\mathbf{P}^{(t)}_{\text{raw}}
                                  + (1-\alpha)\,\mathbf{P}^{(t-1)}_{\text{smooth}},
\qquad \alpha = 0.65$$

Applied independently to $x$ and $y$. The name is \textbf{exponential moving average} (EMA), or
single-pole low-pass filter.

\textbf{Why it works.} Expand the recurrence backwards:

\begin{align*}
\mathbf{P}^{(t)} &= \alpha \mathbf{R}^{(t)} + (1-\alpha)\Big[\alpha \mathbf{R}^{(t-1)}
                    + (1-\alpha)\mathbf{P}^{(t-2)}\Big] \\
                 &= \alpha \mathbf{R}^{(t)} + \alpha(1-\alpha)\mathbf{R}^{(t-1)}
                    + (1-\alpha)^2 \mathbf{P}^{(t-2)} \\
                 &= \alpha \sum_{k=0}^{n} (1-\alpha)^k \mathbf{R}^{(t-k)}
                    + (1-\alpha)^{n+1}\mathbf{P}^{(t-n-1)}
\end{align*}

The current position is a weighted sum of \emph{all} past measurements, with geometrically decaying
weights $\alpha, \alpha(1-\alpha), \alpha(1-\alpha)^2, \dots$: recent frames matter most, old frames
fade out but never vanish entirely.

\textbf{The noise-reduction price.} The weights sum to $1$ in the limit, so a constant true position
is reproduced exactly. Random noise, however, is averaged down: the effective averaging window has
length roughly $1/\alpha \approx 1.5$ frames of ``memory mass'' concentrated in recent frames, and
the residual jitter is reduced by roughly $\sqrt{\alpha/(2-\alpha)} \approx 0.69$ for
$\alpha = 0.65$.

\textbf{The responsiveness price.} If the true position jumps (you flick your finger), the displayed
value approaches it geometrically and never exactly arrives. After one frame it has covered
$\alpha = 65\%$ of the gap; after two, $1-(1-\alpha)^2 = 87.75\%$; after three, $95.7\%$. Visible
lag at 30 FPS is therefore on the order of one to two frames --- about 30--60 milliseconds, below the
threshold at which a user perceives the marker as ``slow''.
\end{mathdeepdive}

\begin{mathdeepdive}{Worked example, and the choice of $\alpha$}
Suppose the smoothed value is $P^{(t-1)} = 300$ and the new raw measurement is $R^{(t)} = 340$.

$$P^{(t)} = 0.65 \times 340 + 0.35 \times 300 = 221 + 105 = 326$$

So the dot moves $26$ pixels of the $40$-pixel gap --- most of the way, but not all of it. The
remaining $14$ pixels will be consumed over the following frames.

Why $\alpha = 0.65$ rather than something else?

\begin{itemize}
  \item $\alpha \to 1$ (say $0.95$): almost no smoothing. Jitter is preserved; the dot is a raw
        measurement with a hat on.
  \item $\alpha \to 0$ (say $0.1$): heavy smoothing. The dot becomes a lazy, laggy ghost that
        visibly trails the finger.
  \item $\alpha = 0.65$: covers $65\%$ of any gap per frame --- fast enough to feel instantaneous at
        video rates, slow enough to flatten single-frame noise.
\end{itemize}

\textbf{First sighting.} On the very first frame there is no previous value, so the recurrence has no
starting point. Blending against zero would make the dot animate in from the top-left corner, which
looks like a bug. Instead, the implementation treats ``no previous'' as a special case and passes the
measurement through unchanged:

\begin{lstlisting}
def ema(previous, current, alpha=0.65):
    if previous is None:
        return (current[0], current[1])
    return (
        alpha * current[0] + (1 - alpha) * previous[0],
        alpha * current[1] + (1 - alpha) * previous[1],
    )
\end{lstlisting}
\end{mathdeepdive}

\begin{edgetrap}{Smoothing state is keyed by hand, not by list index}
MediaPipe does not guarantee the order of its hand list between frames. If your right hand leaves the
frame, or the two hands cross, the hand at index 0 in one frame may be the hand at index 1 in the
next.

If the smoothing history were stored per \emph{index} --- \code{\_smooth\_tip[0]},
\code{\_smooth\_tip[1]} --- a reorder would blend the position of one physical hand with the history
of the other. The visible symptom is a sudden teleport of the marker, which is precisely the
behaviour the filter exists to prevent.

The project keys the state by \textbf{handedness label} instead, falling back to the index only when
the label is unavailable:

\begin{lstlisting}
@staticmethod
def _state_key(hand_index, label):
    return label if label in ("Left", "Right") else f"index:{hand_index}"
\end{lstlisting}

This is not perfect --- handedness itself can flicker --- but it is meaningfully more stable than a
raw index, and it costs nothing. The same key is reused for the gesture debouncer, so a reorder
cannot contaminate gesture history either.
\end{edgetrap}

\section{Gestures and the debouncer}
\label{sec:gesture}

\subsection{From five booleans to a name}

The finger logic produces five booleans in a fixed order --- \code{[thumb, index, middle, ring,
pinky]}. A gesture is simply a recognised pattern in those five bits.

\begin{center}
\rowcolors{2}{slatebg}{white}
\begin{tabular}{@{}l c c c c c l@{}}
\toprule
\rowcolor{white}
\textbf{Gesture} & \textbf{Thumb} & \textbf{Index} & \textbf{Middle} & \textbf{Ring} & \textbf{Pinky}
 & \textbf{Notes} \\
\midrule
Open Palm  & \checkmark & \checkmark & \checkmark & \checkmark & \checkmark & all five extended \\
Fist       &            &            &            &            &            & all five folded \\
Thumbs Up  & \checkmark &            &            &            &            & thumb alone \\
Peace      &            & \checkmark & \checkmark &            &            & the V \\
Rock On    & \checkmark & \checkmark &            &            & \checkmark & index, pinky, thumb \\
(anything else) &   ---   &    ---     &    ---     &    ---     &    ---     & reported as \code{---} \\
\bottomrule
\end{tabular}
\end{center}

The order of the checks in the code matters, and is written so that the more specific patterns are
tested first. Note that ``Open Palm'' (all five true) and ``Thumbs Up'' (only the thumb true) cannot
collide, but reading the code top to bottom shows the intent:

\begin{codecard}
\begin{lstlisting}
if thumb and index and middle and ring and pinky:
    return "Open Palm"
if not (thumb or index or middle or ring or pinky):
    return "Fist"
if thumb and not (index or middle or ring or pinky):
    return "Thumbs Up"
if not thumb and index and middle and not (ring or pinky):
    return "Peace"
if thumb and index and not (middle or ring) and pinky:
    return "Rock On"
return NO_GESTURE
\end{lstlisting}
\end{codecard}

The negated conditions (``not ring and not pinky'') are what make each pattern exact rather than a
prefix match. Without them, ``Peace'' would also match a hand with all five fingers up.

\subsection{Why a raw per-frame answer flickers}

\begin{mathdeepdive}{Temporal majority voting}
During a transition --- say from Fist to Peace --- the fingers do not move in perfect synchrony, and
for a few frames the pattern belongs to no gesture at all. Classifying each frame independently
produces a label that flashes through intermediate states. Worse, a single noisy frame in the middle
of a steady pose can produce a one-frame spike to a wrong label.

The fix is a \textbf{rolling majority vote} over a short history. Formally, with window length $n$
and required agreement $k$, at frame $t$:

$$\text{report } g \iff \sum_{i=t-n+1}^{t} \mathbb{1}\!\left[g_i = g\right] \ \ge\ k,
\qquad n = 5,\ k = 4$$

where $\mathbb{1}[\cdot]$ is the indicator function (1 if the condition holds, 0 otherwise). If no
label reaches $k$, the previously reported label is held.

The implementation is a \code{deque} with \code{maxlen=5}, which gives the sliding window for free,
and a \code{count} on the newest label:

\begin{lstlisting}
def update(self, gesture):
    self._history.append(gesture)
    if self._history.count(gesture) >= self.min_agree:
        self._stable = gesture
    return self._stable
\end{lstlisting}

\textbf{Worked trace.} The window is the last five entries; entries are listed oldest first.

\begin{center}
\rowcolors{2}{slatebg}{white}
\begin{tabular}{@{}r l l l@{}}
\toprule
\rowcolor{white}
\textbf{Frame} & \textbf{Input} & \textbf{Window after append} & \textbf{Returned} \\
\midrule
1 & Peace & P            & \code{---} \\
2 & Peace & P P          & \code{---} \\
3 & Peace & P P P        & \code{---} \\
4 & Peace & P P P P      & \textbf{Peace} \\
5 & Peace & P P P P P    & Peace \\
6 & Fist  & P P P P F    & Peace \\
7 & Fist  & P P P F F    & Peace \\
8 & Fist  & P P F F F    & Peace \\
9 & Fist  & P F F F F    & \textbf{Fist} \\
\bottomrule
\end{tabular}
\end{center}

Reading it carefully: a change of gesture takes exactly four frames to be accepted, and a two-frame
blip can never break through, because two frames cannot reach a count of four in a five-slot window.
That is the precise trade: $\approx 130$\,ms of latency at 30 FPS, in exchange for the complete
elimination of single- and double-frame flicker.
\end{mathdeepdive}

\begin{keyidea}
\textbf{Warm-up behaviour is defined, not accidental.} Before any majority has formed, the stable
label is the placeholder \code{---}, so the HUD shows no gesture for the first three frames of a
session rather than guessing. The alternative --- returning the first raw label immediately --- would
make the debouncer pointless on startup, since the first frame is exactly the one most likely to be
noisy.
\end{keyidea}

\section{The complete decision path}

Putting the chapter together, one frame produces this chain, entirely inside \code{hand\_logic.py}:

\begin{center}
\begin{tikzpicture}[
  node distance=0.32cm,
  st/.style={draw=slateline, fill=slatebg, rounded corners=1mm, align=center,
             inner sep=2.2mm, font=\scriptsize\sffamily, text=navy},
  ar/.style={-{Stealth[length=1.8mm]}, draw=tealx, line width=0.8pt}
]
  \node[st] (a) {normalized\\landmarks};
  \node[st, right=of a] (b) {pixel landmarks\\\code{to_pixel_landmarks}};
  \node[st, right=of b] (c) {palm normal\\\code{palm_normal_z}};
  \node[st, right=of c] (d) {5 booleans\\\code{fingers_state}};
  \node[st, right=of d] (e) {gesture name\\\code{classify_gesture}};
  \node[st, right=of e] (f) {stable label\\\code{GestureDebouncer}};

  \draw[ar] (a) -- (b);
  \draw[ar] (b) -- (c);
  \draw[ar] (c) -- (d);
  \draw[ar] (d) -- (e);
  \draw[ar] (e) -- (f);
\end{tikzpicture}
\captionof{figure}{The per-frame decision chain in the pure-logic module. Every stage is a plain
function of its inputs --- no global state, no camera, no I/O --- which is why the whole chain is
unit-testable with hand-written coordinates.}
\end{center}

\begin{keyidea}
\textbf{Chapter summary.} On a screen $y$ grows downward, so a fingertip is extended when
$y_{\text{tip}} < y_{\text{pip}}$ for the four long fingers. The thumb uses a distance comparison
against the pinky knuckle, which is provably invariant under image-plane rotation and independent of
handedness --- with the adducted ``knife hand'' documented as its known failure. The palm normal's
sign reveals which side of the hand faces the camera. An exponential moving average with
$\alpha = 0.65$ removes jitter at the cost of roughly a frame of latency, and a five-frame,
four-vote rolling majority removes gesture flicker at the cost of roughly four frames of latency.
\end{keyidea}
